\documentclass[12pt]{article}

\usepackage[T1]{fontenc}
\usepackage[utf8]{inputenc}
\usepackage{lmodern}
\usepackage[margin=1in]{geometry}
\usepackage{amsmath,amssymb}
\usepackage{siunitx}
\usepackage{microtype}
\usepackage{booktabs}
\usepackage{enumitem}
\usepackage[hidelinks]{hyperref}

\sisetup{
  detect-all,
  separate-uncertainty=true,
  range-phrase=--,
  range-units=single
}

\title{Falcon 9 Seismo-acoustic Analysis: Results Draft}
\author{}
\date{Draft reconstructed from the canonical analysis outputs}

\begin{document}
\maketitle

\section{Results}

\subsection{Candidate-event catalogue and temporal evolution}

Manual review yielded 384 accepted infrasound picks on HD1--HD3. Correction for
the predicted differential travel times reduced the normal approximately
\SIrange{80}{90}{\milli\second} traversal time across the array and allowed the
picks to be associated using a \SI{0.040}{\second} reduced-time window. The
association produced 154 candidate events, of which 63 contained picks on all
three infrasound channels and 91 contained picks on two channels. Thirteen
accepted picks remained unassociated.

The candidate catalogue extended from 13:07:15.927 to 13:33:54.396 UTC, a span
of approximately \SI{26.64}{\minute}. The activity was strongly episodic rather
than uniformly distributed through this interval (Figure~2). Using the four
display intervals defined in Figure~2, 38 candidates occurred during Phase~I,
3 during Phase~II, 87 during Phase~III, and 6 during Phase~IV; the remaining 20
candidates fell outside the marked intervals. The opening sequence contained the
largest-amplitude signals, whereas Phase~III contained the greatest number of
catalogued impulses. A period of sustained seismic tremor followed the opening
sequence between approximately 13:08:48 and 13:12:00 UTC, with no comparably
continuous signal on the infrasound channels.

The positive-compression/negative-rarefaction algorithm returned a valid pair
for 153 of the 154 candidate events. The sole exception was the principal
explosion, whose waveform was too complex to be represented by a single short
positive--negative cycle. On the initial descriptive classification, 46 events
were labelled robust impulsive events, 82 probable events, 25 weak coherent
candidates, and one---the principal explosion---a complex or noncanonical
waveform. These labels describe waveform quality rather than distinct physical
source types.

Of the 154 candidates, 148 passed the permissive criteria for planar array
analysis, including the explicitly retained principal explosion. Seventy-six
events passed the more conservative peak-to-peak pressure and SNR thresholds and
formed the core planar set. Inspection of the underlying amplitude, SNR, and
coherence distributions did not reveal sharp discontinuities at either set of
thresholds; the core set is therefore a conservative subset of a continuous
quality distribution rather than a separately resolved event population.

\subsection{Opening sequence and key accident stages}

The opening sequence contained four geophysically distinct signals associated
with stages identified in the video chronology: the initial upper-stage failure,
the principal lower-stage explosion, a capsule impact, and a subsequent
capsule-related explosion (Table~2 and Figure~S1). The source-time estimates were
13:07:12.080, 13:07:15.750, 13:07:24.600, and 13:07:25.150 UTC, respectively.

The initial upper-stage failure produced a median-stack positive pressure of
\SI{23}{\pascal}, a negative pressure of \SI{-21}{\pascal}, and a peak-to-peak
pressure of \SI{44}{\pascal}. Its three-component vector PGV was
$7.85\times10^{-5}\ \si{\metre\per\second}$. The best planar solution gave a
back azimuth of \ang{201.5} and an apparent speed of
\SI{355.5}{\metre\per\second}.

Approximately \SI{3.67}{\second} later, the principal explosion produced the
largest acoustic and seismic amplitudes in the catalogue (Figure~3). Using the
aligned median-stack estimator applied consistently to all events, the positive
pressure was \SI{1353}{\pascal}, the negative pressure was
\SI{-84}{\pascal}, and the peak-to-peak pressure was \SI{1437}{\pascal}. The
positive pressure was approximately 59 times that of the initial upper-stage
signal. Assuming geometric spreading proportional to $1/r$, the stack-positive
pressure corresponds to \SI{1919}{\pascal} at \SI{1}{\kilo\metre}. Its peak
sound-pressure level at BCHH was \SI{156.6}{\decibel} re
\SI{20}{\micro\pascal}.

Measurements made separately on the three infrasound sensors gave positive
peaks of 1.335, 1.546, and \SI{1.423}{\kilo\pascal} on HD1, HD2, and HD3,
respectively. The corresponding negative extrema were $-193$, $-231$, and
\SI{-234}{\pascal}. Thus, the median of the sensor-specific positive extrema was
\SI{1.423}{\kilo\pascal}, equivalent to approximately
\SI{2.014}{\kilo\pascal} at \SI{1}{\kilo\metre}. The difference from the
median-stack result reflects the use of two distinct estimators: the stack is
evaluated sample by sample after alignment, whereas the sensor-specific extrema
need not occur at exactly the same relative time.

The principal pressure waveform comprised an abrupt positive compression
followed by a longer, irregular decay and substantially smaller rarefaction. It
therefore did not display the compact, approximately symmetric positive--negative
morphology of many weaker catalogue events. Its shape qualitatively resembles a
blast pulse with an extended decay, although no parametric Friedlander-waveform
fit was performed. The seismic response persisted considerably longer than the
initial acoustic compression and contained large-amplitude motion on all three
components. The three-component vector PGV reached
$3.57\times10^{-3}\ \si{\metre\per\second}$, the largest value measured in the
catalogue.

The two capsule-related signals were separated by approximately
\SI{0.55}{\second}. The first produced stack pressures of $+276$ and
\SI{-102}{\pascal}, a peak-to-peak pressure of \SI{378}{\pascal}, and vector
PGV of $6.98\times10^{-4}\ \si{\metre\per\second}$. The second produced
$+291$ and \SI{-113}{\pascal}, a peak-to-peak pressure of
\SI{404}{\pascal}, and vector PGV of
$9.02\times10^{-4}\ \si{\metre\per\second}$. Their planar back azimuths were
\ang{201.5} and \ang{201.3}, and their apparent speeds were 351.5 and
\SI{351.0}{\metre\per\second}, respectively. Both directions and speeds were
consistent with propagation from the launch-complex sector.

\subsection{Planar-wave array results}

The 148 analyzed events produced a narrow distribution of best-fitting source
directions (Figure~5). The median planar back azimuth was \ang{200.4}, with an
interquartile range (IQR) of \ang{1.3}. The complete analyzed range was
\angrange{195.8}{208.2}, although the extremes were dominated by weak or broad
coherence surfaces. Within the 76-event core set, the median was \ang{200.5},
the IQR was \ang{1.2}, and the range was \angrange{196.8}{205.3}. The core
median lay approximately \ang{1.1} clockwise from the surveyed direction to
SLC-40 (\ang{199.4}). At a range of \SI{1.42}{\kilo\metre}, an angular
difference of \ang{1.1} corresponds to approximately \SI{27}{\metre} in the
cross-range direction, but this offset is not uniquely attributable to source
position because geometry, atmospheric structure, waveform differences, and
finite-range effects can produce comparable shifts.

The median apparent speed for all analyzed events was
\SI{352.0}{\metre\per\second}, with an IQR of
\SI{12.6}{\metre\per\second}. The core set was more tightly concentrated, with
a median of \SI{351.25}{\metre\per\second} and an IQR of
\SI{6.75}{\metre\per\second}. Both medians closely matched the
meteorologically estimated effective speed of
\SI{350.98}{\metre\per\second}. The wider complete-catalogue range of
\SIrange{280}{440}{\metre\per\second} included solutions that reached the
search boundaries and was associated mainly with weak or poorly constrained
events.

The median joint coherence score was 0.819 for all analyzed events and 0.942 for
the core set. Fifty planar solutions met the high-quality coherence and
window-stability criteria, 46 were classified as probable, and 52 were weak or
broad. Among the core events, 47 were high quality, 21 probable, and only 8 weak
or broad. Perturbing the scoring windows produced median standard deviations of
\ang{0.29} in back azimuth and \SI{1.89}{\metre\per\second} in speed across all
analyzed events. For the core set these decreased to \ang{0.24} and
\SI{1.44}{\metre\per\second}, respectively.

The pairwise results supported the visual observation that HD2 and HD3 were
usually more similar to one another than either was to HD1. Across all analyzed
events, the median HD2--HD3 correlation was 0.928, compared with 0.645 for
HD1--HD2 and 0.644 for HD1--HD3. Within the core set, the corresponding medians
were 0.973, 0.886, and 0.863. The median HD1 correlation deficit---the HD2--HD3
correlation minus the mean of the two HD1 pair correlations---was 0.197 for all
analyzed events and 0.068 for the core set. HD1 therefore reduced coherence most
strongly among the weaker events, while remaining well correlated for most core
events.

The principal explosion gave a planar back azimuth of \ang{200.2}, an apparent
speed of \SI{352.5}{\metre\per\second}, and a coherence score of 0.9989. Its
solution was exceptionally stable under changes to the scoring window. The
composite coherence surface for the core catalogue also formed a compact maximum
near the surveyed launch-complex direction and the meteorologically predicted
speed, demonstrating that the catalogue-scale result was not controlled solely
by the principal explosion.

\subsection{Finite-range circular-wavefront results}

Replacing the plane-wave approximation with exact horizontal distances to a
source fixed at SLC-40 produced only modest changes. The median preferred
circular-wavefront speed was \SI{355.75}{\metre\per\second}, compared with
\SI{352.0}{\metre\per\second} for the planar model. The circular model's median
coherence score was lower than the planar optimum by only 0.0047. This small
difference is expected because the planar model can vary both direction and
speed, whereas the fixed-source calculation constrains the source direction.

For the principal explosion, the fixed-SLC-40 circular solution preferred
\SI{355.0}{\metre\per\second} and gave a coherence score of 0.9965, compared
with \SI{352.5}{\metre\per\second} and 0.9989 for the planar solution. Thus,
wavefront curvature across the approximately \SI{30}{\metre} aperture did not
materially alter the inferred propagation characteristics.

The range-profile analysis demonstrated that the array could not determine
source range. At a coherence-score reduction of 0.01, the along-range profile
remained open at the search boundary for all 148 events, and 66 event maxima lay
on an along-range boundary. Cross-range behavior was better constrained but
remained broad: only 7 profiles were open at the cross-range boundaries, the
median best cross-range offset was \SI{25}{\metre}, and the median full width at
a score reduction of 0.01 was approximately \SI{115}{\metre}. The apparent
clustering of trial locations away from SLC-40 therefore does not constitute a
resolved source-location result. It primarily reflects a broad trade-off among
cross-range position, propagation speed, and small systematic delay errors.

\subsection{Event amplitudes through time}

Across all 154 candidates, median stack peak-to-peak pressure was
\SI{29.5}{\pascal}, with an interquartile range of
\SIrange{16.2}{51.3}{\pascal}. Median vector PGV was
$4.26\times10^{-5}\ \si{\metre\per\second}$, with an interquartile range of
$2.27\times10^{-5}$ to
$1.05\times10^{-4}\ \si{\metre\per\second}$. Both acoustic and seismic
amplitudes were largest during the opening sequence, declined rapidly after the
principal explosion and capsule-related signals, and increased again during the
dense Phase~III sequence (Figure~6). Later events were generally weaker than the
opening explosions, although several isolated impulses remained measurable late
in the record.

The principal explosion was the largest event by both pressure and vector PGV.
The next two candidates immediately following it reached peak-to-peak pressures
of approximately 699 and \SI{632}{\pascal} and vector PGVs of approximately
$2.78\times10^{-3}\ \si{\metre\per\second}$. This concentration of the largest
signals within a few seconds shows that the opening failure sequence dominated
the measured acoustic and seismic amplitudes, even though most catalogued events
occurred later.

The median dimensional pressure-to-PGV ratio for the full catalogue was
$5.70\times10^{5}\ \si{\pascal\per\metre\second}$. The ratio showed much more
scatter among events with vector PGV below approximately
$10^{-4}\ \si{\metre\per\second}$. This behavior coincided with decreasing
seismic SNR and reduced capture of the event maximum, indicating that at least
part of the increased low-amplitude scatter resulted from the seismic noise
floor rather than a demonstrable change in acoustic--seismic coupling.

\subsection{Acoustic pressure versus ground velocity}

Twenty-seven events met the joint acoustic-SNR, seismic-SNR, and window-capture
criteria for regression (Figure~7). Within this subset, peak-to-peak pressure and
vector PGV were strongly correlated in logarithmic space. The fitted power-law
exponent was
\begin{equation}
  b = 0.868 \pm 0.048,
\end{equation}
where the quoted uncertainty is the ordinary least-squares standard error. The
95\% event-bootstrap interval was 0.734--0.971. The regression explained 93.0\%
of the variance in $\log_{10}$ pressure ($R^2=0.930$).

The fitted exponent was slightly below unity, implying a modestly sublinear
relationship over the observed amplitude range. Nevertheless, the proportional
model
\begin{equation}
  p_{\mathrm{p2p}} = K\,\mathrm{PGV}
\end{equation}
with
\begin{equation}
  K = \SI{0.447}{\mega\pascal\per\metre\second}
\end{equation}
provided a nearly comparable description. The free power-law and proportional
models had base-10 logarithmic RMSE values of 0.128 and 0.146, respectively,
equivalent to typical multiplicative residual factors of approximately 1.34 and
1.40. The small difference in predictive scatter indicates that a nearly
constant pressure-to-PGV ratio remains a useful first-order description, even
though the unconstrained regression favored a sublinear exponent.

The principal explosion lay within the high-quality regression subset. Its
pressure-to-PGV ratio was approximately
$4.02\times10^{5}\ \si{\pascal\per\metre\second}$, comparable to the fitted
proportionality constant. The largest events therefore extended the general
acoustic--seismic amplitude trend rather than forming a clearly separate
population. In contrast, many of the weakest candidates departed substantially
from the fitted trend and should not be used to infer changes in energy
partitioning without a more explicit treatment of detection limits and
uncertainties in both measurements.

\clearpage
\section*{Results items requiring confirmation before submission}

\begin{enumerate}[leftmargin=*]
  \item \textbf{Principal pressure reported in the abstract and conclusions.}
  Decide whether the headline value will be the catalogue-consistent aligned
  stack positive pressure (\SI{1.353}{\kilo\pascal};
  \SI{1.919}{\kilo\pascal} at \SI{1}{\kilo\metre}) or the median of the three
  sensor-specific positive extrema (\SI{1.423}{\kilo\pascal};
  \SI{2.014}{\kilo\pascal} at \SI{1}{\kilo\metre}). Both are defensible but
  represent different estimators.

  \item \textbf{Phase definitions.} Confirm that the four intervals drawn in
  Figure~2 are intended as formal phases rather than visual annotations. If they
  are interpretive, their physical definitions should be stated before reporting
  counts per phase.

  \item \textbf{Key-event terminology.} The canonical key-event file still uses
  legacy identifiers \texttt{capsule\_pulse\_1} and
  \texttt{capsule\_pulse\_2}. Confirm the interpretations ``capsule impact'' and
  ``capsule explosion'' from the video analysis.

  \item \textbf{Source-direction offset.} The approximately \ang{1} median
  offset from the surveyed SLC-40 direction should not be interpreted as source
  displacement until the geometry convention, atmospheric crosswind, and
  channel timing have been fully audited.

  \item \textbf{Circular-wavefront source positions.} Do not report the best
  cross-range grid points as event locations. The broad profiles and unresolved
  along-range dimension support only a resolution-limit result.

  \item \textbf{Regression model.} The bootstrap interval excludes a unit
  exponent, but the two model RMSE values differ only slightly. Decide whether
  the paper should emphasize the statistically sublinear fit, the practical
  adequacy of a constant ratio, or both.

  \item \textbf{Energy results.} Acoustic and seismic energy and equivalent
  magnitude values are omitted because the current canonical notebooks do not
  regenerate them.
\end{enumerate}

\end{document}
